\section{SWE-bench Trace Audit: Registry vs. Repo Workflow}
\label{app:swebench-trace-audit}

HAL Generalist and SWE-Agent both expose code-generation, code-execution,
shell, file-read, and file-edit affordances on SWE-bench. In HAL Generalist,
these affordances pass through a smolagents-style code-agent registry with
tools such as \texttt{python\_interpreter}, \texttt{execute\_bash},
\texttt{inspect\_file\_as\_text}, \texttt{edit\_file},
\texttt{file\_content\_search}, and \texttt{final\_answer}. SWE-Agent exposes
a repo-repair workflow with a specialized action protocol, persistent editor
semantics, file-map context, diff/review state, and an explicit submit action.
The audit below therefore focuses on the model-facing workflow difference,
rather than treating live execution or editing as the distinguishing factor.

To make this distinction concrete, we audited matched GPT-5 traces on
SWE-bench Verified Mini. The HAL Generalist run
\texttt{\detokenize{swebench_verified_mini_hal_generalist_gpt520250807_1755463923}}
and the SWE-Agent run
\texttt{\detokenize{swebench_verified_mini_sweagentgpt520250807_1754592641}}
cover the same 50 tasks. HAL Generalist fails while SWE-Agent succeeds on 26
of these tasks. The joinable extraction artifacts are stored in
\texttt{irt\_data/trace\_analysis/swebench\_hal\_vs\_sweagent/}: the full
candidate set is \texttt{gpt5\_hal\_fail\_swe\_success\_tasks.csv}, and the
paired trace examples below are \texttt{gpt5\_trace\_examples.csv}.

\begin{table}[h]
\centering
\footnotesize
\setlength{\tabcolsep}{2pt}
\begin{tabular}{>{\raggedright\arraybackslash}p{0.15\linewidth}>{\raggedright\arraybackslash}p{0.24\linewidth}>{\raggedright\arraybackslash}p{0.29\linewidth}>{\raggedright\arraybackslash}p{0.25\linewidth}}
\toprule
Task & Issue & HAL Generalist failed trace & SWE-Agent succeeded trace \\
\midrule
Django 12143 &
Regex escaping in Django admin formset prefixes &
8 LLM calls; uses registry tools including \texttt{file\_content\_search},
\texttt{edit\_file}, and \texttt{execute\_bash}; final message says the
referenced upstream code is absent and no patch is generated. &
25 LLM calls; uses \texttt{bash}, \texttt{str\_replace\_editor}, and
\texttt{submit}; final message reports locating the ModelAdmin edited-PK
helper, applying \texttt{re.escape}, and verifying with a reproduction script. \\
Django 12050 &
Preserve list-vs-tuple type in query lookup resolution &
19 LLM calls; uses search, edit, bash, and a Python interpreter; final message
is still in patch-construction mode and the submitted result fails grading. &
16 LLM calls; uses repository shell/editor actions and submit; trace records a
plan to inspect ORM code, reproduce the coercion bug, patch it, and rerun the
reproduction script. \\
Sphinx 9698 &
Remove parentheses from Sphinx Python-property index entries &
19 LLM calls; uses repeated registry searches and edits; final message returns
a patch description, but the benchmark marks it unresolved. &
42 LLM calls; uses repeated shell/editor operations and submit; final message
reports a minimal change in the Sphinx Python-domain implementation. \\
\bottomrule
\end{tabular}
\caption{Matched GPT-5 examples where HAL Generalist fails and SWE-Agent
succeeds on the same SWE-bench Verified Mini task. Both harnesses can execute
and edit code; the visible difference is the model-facing workflow: a broad
registry of general tools versus a repo-specific repair protocol.}
\label{tab:swebench-trace-examples}
\end{table}

The examples suggest a concrete follow-up ablation: keep the same base model
and sandbox budget, but vary only the repo-repair interface. For HAL
Generalist, this could mean adding a SWE-Agent-like file-map/edit/review/submit
layer while preserving its underlying execution tools. For SWE-Agent, this
could mean disabling file-map or diff-state context while leaving shell and
editor access intact. Such paired ablations would separate the value of live
execution from the value of repo-specialized workflow structure.
